\chapter{ภาพรวมโครงการวิจัยและคณะผู้วิจัย}
\label{chap:overview}

\section{ความเป็นมาและความสำคัญของปัญหา}
การตรวจวัดสมบัติทางกายภาพและเคมีของดิน โดยเฉพาะอย่างยิ่งค่าความเป็นกรด-ด่างของดิน (Soil pH) มีความสำคัญอย่างยิ่งยวดต่อผลิตภาพและคุณภาพของผลผลิตทางการเกษตร โดยเฉพาะในพื้นที่ภาคตะวันออกของประเทศไทยซึ่งเป็นแหล่งปลูกทุเรียน (\textit{Durio zibethinus} Murr.) ที่สำคัญระดับโลก สภาพดินที่มีความเป็นกรดรุนแรง (pH ต่ำกว่า 5.0) ส่งผลให้ธาตุอาหารหลัก เช่น ฟอสฟอรัส โพแทสเซียม แคลเซียม และแมกนีเซียม ถูกตรึงอยู่ในดินและไม่สามารถดูดซึมผ่านรากพืชได้อย่างมีประสิทธิภาพ ขณะเดียวกันยังปลดปล่อยไอออนของอะลูมิเนียมและแมงกานีสออกมาในระดับที่เป็นพิษต่อรากพืช

แม้ว่าเกษตรกรจะตระหนักถึงความสำคัญของการวัดค่า pH ดิน แต่เครื่องมือวัดภาคสนามที่ใช้งานอยู่ในปัจจุบันกลับมีข้อจำกัดอย่างมาก กล่าวคือ ชุดตรวจวัดแบบเทียบแถบสีหรือสารเคมีมีความคลาดเคลื่อนเชิงบุคคลสูงและไม่สะดวกต่อการบันทึกข้อมูลดิจิทัล ในขณะที่หัววัดเคมีไฟฟ้าแบบอิเล็กทรอนิกส์ราคาประหยัดมักประสบปัญหาความคลาดเคลื่อนจากการเปลี่ยนแปลงของอุณหภูมิภาคสนาม (Temperature Drift) และพฤติกรรมที่ไม่เป็นเชิงเส้น (Non-linearity) ของเยื่อแก้วสัมผัส ทำให้ค่าที่อ่านได้คลาดเคลื่อนสูงถึง $\pm 0.3 - 0.8$ หน่วย pH ซึ่งเกินกว่าเกณฑ์ความคลาดเคลื่อนที่ยอมรับได้ในการเกษตรแม่นยำ

โครงการวิจัยนี้จึงมุ่งเน้นการบูรณาการหลักการแรกทางฟิสิกส์เคมี (First-Principles Physics) เข้ากับปัญญาประดิษฐ์ประมวลผลบนขอบ (Edge AI) เพื่อพัฒนาชุดเครื่องมือวัดอัจฉริยะที่สามารถชดเชยอุณหภูมิและความคลาดเคลื่อนทางกายภาพได้แบบเรียลไทม์บนสมาร์ทโฟนโดยตรง

\section{วัตถุประสงค์ของการวิจัย}
วัตถุประสงค์หลักของการวิจัยมี 4 ประการ ได้แก่
\begin{enumerate}
    \item เพื่อศึกษาและจำลองผลกระทบของอุณหภูมิและความไม่เป็นเชิงเส้นของหัววัดเคมีไฟฟ้าต่อค่าศักย์ไฟฟ้าตามสมการเนิร์นสต์
    \item เพื่อออกแบบและพัฒนาต้นแบบระบบวัดพารามิเตอร์ดินเชื่อมต่อผ่านพอร์ต USB-C มาตรฐาน RS485 Modbus RTU พร้อมระบบบริหารจัดการพอร์ตเพื่อป้องกันปัญหาการแย่งสิทธิ์เชื่อมต่อระหว่างแอปพลิเคชัน
    \item เพื่อพัฒนาและฝึกฝนแบบจำลองปัญญาประดิษฐ์โครงข่ายพหุนามประมวลผลบนขอบ (Pure Dart Edge ML OLS Trainer) ที่สามารถชดเชยความคลาดเคลื่อนได้บนสมาร์ทโฟนโดยไม่พึ่งพาเซิร์ฟเวอร์
    \item เพื่อทดสอบความใช้ได้ของวิธีตามมาตรฐานสากล และประเมินสมรรถนะภาคสนามในแปลงปลูกทุเรียนจริง 30 จุดในเขตจังหวัดจันทบุรีและตราด
\end{enumerate}

\section{ข้อมูลโครงการวิจัยและสัดส่วนการวิจัย}
โครงการวิจัยได้รับการสนับสนุนงบประมาณจากกองทุนวิจัย มหาวิทยาลัยราชภัฏรำไพพรรณี ประจำปีงบประมาณ พ.ศ. 2569 โดยมีคณะผู้วิจัยและสัดส่วนภาระงานวิจัย ดังแสดงในตารางที่ \ref{tab:researchers}

\begin{table}[H]
\caption{คณะผู้วิจัยและสัดส่วนการวิจัย โครงการพัฒนาเครื่องวัด pH ดิน SoilpHTxAI}
\label{tab:researchers}
\centering
\begin{tabularx}{\textwidth}{p{4.2cm}p{3.5cm}cZ}
\toprule
\textbf{ชื่อ-สกุล} & \textbf{ตำแหน่งในโครงการ} & \textbf{สัดส่วน} & \textbf{สังกัดหน่วยงาน} \\
\midrule
อาจารย์ธนพัฒน์ ถิระวุฒิ & หัวหน้าโครงการวิจัย & 50\% & หลักสูตร คบ.ฟิสิกส์ มรภ.รำไพพรรณี \\
ผศ.ดร.ชีวะ ทัศนา & นักวิจัย & 30\% & หน่วยวิจัย AI เกษตรดิจิทัล มรภ.รำไพพรรณี \\
รศ.ดร.นันทพร มูลรังษี & นักวิจัย & 10\% & สาขาวิชาเคมี มรภ.รำไพพรรณี \\
รศ.ดร.นิภัทร เปี่ยมอรุณ & นักวิจัย & 10\% & คณะวิทยาศาสตร์และเทคโนโลยี \\
\bottomrule
\end{tabularx}
\end{table}

\begin{academicbox}{หน่วยวิจัยปัญญาประดิษฐ์เพื่อเกษตรดิจิทัล (AI for Digital Agriculture Unit)}
ระบบฮาร์ดแวร์ ซอฟต์แวร์แอปพลิเคชัน และเอกสารคู่มือฉบับสมบูรณ์นี้ จัดทำขึ้นโดย ผู้ช่วยศาสตราจารย์ ดร.ชีวะ ทัศนา ร่วมกับคณะผู้วิจัย หน่วยวิจัยปัญญาประดิษฐ์เพื่อเกษตรดิจิทัล คณะวิทยาศาสตร์และเทคโนโลยี มหาวิทยาลัยราชภัฏรำไพพรรณี เพื่อใช้เป็นต้นแบบนวัตกรรมเกษตรอัจฉริยะที่ถ่ายทอดสู่ชุมชนและเกษตรกรในพื้นที่ภาคตะวันออก
\end{academicbox}

\section{กรอบแนวคิดและสถาปัตยกรรมรวมของระบบ}
ระบบ SoilpHTxAI ประกอบด้วย 4 องค์ประกอบสำคัญ ได้แก่ โพรบเซนเซอร์วัดดินมาตรฐานอุตสาหกรรม, ตัวแปลงสัญญาณ USB-C Serial, แอปพลิเคชันสมาร์ทโฟนที่ผสาน Pure Dart Edge ML และระบบกล้องถ่ายภาพดินพร้อม Geotagging ดังแสดงในภาพที่ \ref{fig:system_overview}

\begin{figure}[H]
\centering
\resizebox{0.88\textwidth}{!}{
\begin{tikzpicture}[
    node distance=1.8cm,
    block/.style={rectangle, draw=rbruNavy, fill=white, text width=3.8cm, text centered, rounded corners=6pt, minimum height=1.3cm, line width=1pt, drop shadow={opacity=0.2}},
    ai_block/.style={rectangle, draw=rbruEmerald, fill=rbruEmerald!10, text width=3.8cm, text centered, rounded corners=6pt, minimum height=1.3cm, line width=1.5pt, drop shadow={opacity=0.3}},
    line/.style={draw=rbruNavy, -{Latex[length=3mm]}, line width=1.5pt}
]
    \node [block] (sensor) {\textbf{โพรบเซนเซอร์ดิน 4-in-1}\\pH, อุณหภูมิ, ความชื้น, EC};
    \node [block, right=1.6cm of sensor] (modbus) {\textbf{วงจรสื่อสารแปลงสัญญาณ}\\RS485 to USB-C Serial\\(Baud Rate 9600 bps)};
    \node [ai_block, below=1.4cm of modbus] (app) {\textbf{แอปพลิเคชัน SoilpHTxAI}\\การชดเชยอุณหภูมิ Edge AI\\(Pure Dart OLS Trainer)};
    \node [block, left=1.6cm of app] (camera) {\textbf{กล้อง Smart Soil Vision}\\Telemetry HUD Overlay\\GPS Geotag ละติจูด-ลองจิจูด};
    \node [block, below=1.4cm of app] (cloud) {\textbf{Cloud \& Drive Sync}\\ส่งออกไฟล์ .CSV / .JSON\\Google Drive / REST API};

    \path [line] (sensor) -- node[above]{\small RS485 Modbus} (modbus);
    \path [line] (modbus) -- node[right]{\small USB-C OTG} (app);
    \path [line] (camera) -- node[above]{\small ภาพถ่าย+พิกัด} (app);
    \path [line] (app) -- node[right]{\small ส่งออกข้อมูล} (cloud);
\end{tikzpicture}
}
\caption{แผนภาพแสดงสถาปัตยกรรมรวมของระบบวัดและวิเคราะห์ pH ดิน SoilpHTxAI}
\label{fig:system_overview}
\end{figure}

\clearpage
